\documentclass[11pt,a4paper]{article}
\usepackage[T1]{fontenc}
\usepackage[utf8]{inputenc}
\usepackage{times}
\usepackage[margin=2.4cm]{geometry}
\usepackage{booktabs}
\usepackage{siunitx}
\usepackage{amsmath}
\usepackage{graphicx}
\usepackage[table]{xcolor}
\usepackage[hidelinks]{hyperref}
\usepackage[numbers,sort&compress]{natbib}
\usepackage{titlesec}
\usepackage{enumitem}
\usepackage{caption}

\definecolor{nile}{HTML}{0B3C5D}
\definecolor{evap}{HTML}{D96F32}
\definecolor{rule}{HTML}{BFCBD3}
\titleformat{\section}{\normalfont\Large\bfseries\color{nile}}{\thesection}{0.6em}{}
\titleformat{\subsection}{\normalfont\large\bfseries\color{nile}}{\thesubsection}{0.6em}{}
\captionsetup{font=small,labelfont={bf,color=nile}}
\sisetup{group-separator={,},detect-weight=true}
\setlist{itemsep=2pt,parsep=0pt,topsep=4pt}

\newcommand{\keyfinding}[1]{%
  \begin{center}\fcolorbox{rule}{black!3}{%
  \begin{minipage}{0.94\linewidth}\vspace{4pt}#1\vspace{4pt}\end{minipage}}\end{center}}

\title{\color{nile}\textbf{Same Budget, More Water}\\[2mm]
\large\normalfont Turning the WaPOR evaporation--transpiration split into a
ranked allocation of Egypt's on-farm modernisation budget}
\author{WaPOR Hackathon 2026 --- technical documentation}
\date{19 August 2026}

\begin{document}
\maketitle

\begin{abstract}
\noindent
FAO's WaPOR v3 database reports evaporation ($E$) and transpiration ($T$) as
separate layers \citep{fao_wapor_v3}. Transpiration passed through a stoma and
grew a crop; evaporation left the field having grown nothing. This work converts
that distinction into an ordering over spending: for each administrative unit in
Egypt's Nile Delta and Valley it estimates the volume of non-beneficial
evaporation that is recoverable by low-cost on-farm measures, the annualised cost
of recovering it, and hence a cost per cubic metre. Spending a fixed budget down
that ordering recovers \SI{25}{\percent} more water than distributing the same
budget in proportion to irrigated area. The work is designed to sit immediately
downstream of IPAT, the irrigation assessment tool already operating on the
Egyptian Ministry of Water Resources and Irrigation platform \citep{fao_ipat},
which reports \emph{where} performance is poor but not \emph{where to spend}.
Two findings contradicted the authors' prior hypothesis and are reported in full,
together with a self-audit of how much of the result is attributable to the
satellite rather than to coarse agronomic knowledge. A substantial methodological
objection --- that WaPOR's $E$/$T$ separation may be insufficiently accurate for
this use \citep{chukalla2022} --- is stated and addressed rather than omitted.
\end{abstract}

\section{The decision, and the gap in the toolchain}

Egypt allocates a fixed share of the Nile among roughly 110 million people, with
agriculture consuming over \SI{85}{\percent} of the national water allocation
across approximately \SI{4}{\percent} of the land area \citep{ayyad2019}. The
Ministry of Water Resources and Irrigation (MWRI) is responsible for that
allocation and, under the National Water Resources Plan 2017--2037, for
rationalising water use and rehabilitating conveyance infrastructure
\citep{nwrp2037}. It spends an annual on-farm modernisation budget and is
accountable for what that budget delivers.

FAO's applications catalogue now lists more than twenty operational WaPOR tools
\citep{fao_catalogue}. Egypt has its own: IPAT is a web tool hosted on the MWRI
platform, built on WaPOR v3 Level 3 (\SI{20}{\metre}) data covering roughly
1.5\,Mha of the Middle and West Delta across sprinkler, drip and
flood systems, co-developed with IHE Delft \citep{fao_ipat}. It reports where
water productivity is high and low.

\keyfinding{The gap is not data, and it is not accuracy. A planning officer with
a fixed budget must decide \emph{which} governorates receive it and in what
proportion. No output in the existing chain is denominated in the units of that
decision. The tool reports millimetres per dekad; the decision is made in
currency, hectares and cubic metres.}

\section{Method}

\subsection{Seasonal aggregation}

WaPOR v3 Level 2 dekadal rasters were summed over the 18 dekads spanning
1 May to 31 October 2024, the summer (\emph{Nili}/\emph{Sefi}) season. Two
corrections are required and neither is optional. WaPOR stores dekadal values as
scaled integers with a scale property of \num{0.1}, and the scaled quantity is a
rate in \si{\milli\metre\per\day}, not a total per dekad. Dekads are 10, 10, and
8--11 days long, so
\begin{equation}
X_{\text{season}} \;=\; \sum_{d=1}^{18} \; \mathrm{DN}_d \times s \times n_d ,
\qquad s = 0.1,\; n_d = \text{days in dekad } d .
\end{equation}
Because $s n_d \approx 1$, the raw integer coincidentally reads as millimetres
per dekad, and a seasonal total computed with \emph{neither} correction is also
superficially plausible. The raw value is therefore not a usable diagnostic; only
the seasonal total distinguishes the two cases.

\subsection{Separating cropland from desert, and from water}

Outside the irrigated Nile system there is effectively no water and no rain, so
seasonal evapotranspiration collapses towards zero. A seasonal threshold of
$\text{AETI} > \SI{200}{\milli\metre}$ therefore separates irrigated land from
desert using the same data already being analysed, with no additional product,
licence or cost.

Open water is a more serious problem, because it evaporates freely all season and
would masquerade as extremely inefficient cropland. Three filters are applied in
combination: JRC Global Surface Water occurrence above \SI{50}{\percent}
\citep{pekel2016}; a beneficial-fraction floor of $T/\text{AETI} > 0.15$, below
which a pixel cannot be a crop; and an absolute ceiling of
$E < \SI{600}{\milli\metre}$ per season. Section~\ref{sec:contamination} explains
why the third filter proved necessary.

\subsection{From evaporation to an ordering over spending}

For each administrative unit $i$ and applicable intervention $j$:
\begin{align}
V_{ij} &= E_i \cdot A_i \cdot 10 \cdot r_j \cdot m_j
  && [\si{\cubic\metre}\,\text{yr}^{-1}] \\
C_{ij} &= \frac{c_j}{L_j} \cdot A_i
  && [\text{USD}\,\text{yr}^{-1}] \\
\kappa_{ij} &= C_{ij} / V_{ij}
  && [\text{USD}\,\text{m}^{-3}]
\end{align}
where $E_i$ is mean seasonal evaporation over masked cropland
(\si{\milli\metre}), $A_i$ the observed irrigated area (\si{\hectare}), $r_j$ the
evaporation-reduction fraction, $m_j$ the number of seasons per year over which
the measure acts, $c_j$ the capital cost per hectare and $L_j$ the service life.
Pairs are sorted by $\kappa$ ascending and funded until the budget is exhausted.
Because a unit may be partially treated, the greedy solution is optimal.

\subsection{Comparator}

The baseline distributes the same budget in proportion to irrigated area, using
the same one-measure-per-unit rule so that the constraint cancels. This is a
\emph{neutral prior} --- what one would do knowing only where the cropland lies
--- and explicitly not a claim about current MWRI practice, for which no citation
is available. Against a literal equal split the reported advantage would be
considerably larger; quoting that number would be comparing against a weaker
counterfactual than honesty permits.

\subsection{Intervention coefficients}

\begin{table}[h]\centering\small
\caption{Intervention model. These are literature-order estimates, not
measurements, and every value is an editable input in the delivered tool.}
\begin{tabular}{lS[table-format=4.0]S[table-format=2.0]S[table-format=1.2]S[table-format=1.0]l}
\toprule
Measure & {USD\,ha$^{-1}$} & {Life, yr} & {$r_j$} & {$m_j$} & Note \\
\midrule
Laser land levelling & 75 & 5 & 0.10 & 2 & Deployed at scale in Egypt \\
Alternate wetting \& drying & 25 & 1 & 0.10 & 1 & Evaporation component only \\
Straw / plastic mulching & 260 & 1 & 0.30 & 1 & Never selected \\
Drip retrofit & 1800 & 10 & 0.45 & 2 & Never selected \\
\bottomrule
\end{tabular}
\label{tab:interventions}
\end{table}

The coefficient for alternate wetting and drying warrants comment. The saving
usually attributed to AWD in the literature is in water \emph{applied} ---
percolation and seepage --- rather than in evapotranspiration
\citep{bouman2001}. In a closed, high-reuse system such as the Nile Delta,
percolation is largely recovered downstream and therefore is not a saving at
basin scale \citep{perry2007,vanhalsema2012}. Only the evaporation component is
counted here, which is why $r_j = 0.10$ rather than the 0.25--0.30 commonly
quoted. Two of the four measures are never selected at any budget between
\$6M and \$80M; they are retained for completeness, not because they do work.

\section{Results}

Observed irrigated area totalled 3.56\,Mha against
3.21\,Mha of reported cultivated area --- an \SI{11}{\percent}
overshoot, in the expected direction for a satellite that also observes
orchards, canal banks and reclamation the statistics omit. As an internal check,
$E + T + I$ summed to AETI exactly, confirming the scale and day-count
corrections.

\begin{table}[h]\centering\small
\caption{Allocation of a \$12M budget. Water figures are annual.}
\begin{tabular}{llS[table-format=1.2]S[table-format=3.1]S[table-format=1.3]}
\toprule
& Governorate & {Funded, \$M} & {Mm$^3$\,yr$^{-1}$} & {USD\,m$^{-3}$} \\
\midrule
1 & Aswan     & 1.19 & 29.5 & 0.040 \\
2 & Qena      & 2.74 & 67.8 & 0.040 \\
3 & Sohag     & 2.17 & 50.8 & 0.043 \\
4 & Assiut    & 2.23 & 48.3 & 0.046 \\
5 & Beni Suef & 1.97 & 41.7 & 0.047 \\
6 & Qalyubia  & 1.12 & 22.9 & 0.049 \\
7 & Minya     & 0.58 & 11.8 & 0.049 \\
\midrule
\multicolumn{2}{l}{\textbf{Targeted}}          & 12.00 & 273.0 & 0.044 \\
\multicolumn{2}{l}{Area-proportional baseline} & 12.00 & 218.1 & 0.055 \\
\bottomrule
\end{tabular}
\label{tab:allocation}
\end{table}

\keyfinding{Ordering the same budget by cost of recovered water yields
\SI{25.2}{\percent} more water than distributing it by irrigated area, and the
same figure on unit cost --- so the gain is not an artefact of spending more.}

\subsection{Equity is priced rather than assumed away}

Strictly cheapest-first funds seven governorates out of eighteen. No ministry
operating under an equity mandate would sign that, so a recommendation which
ignores the constraint is a recommendation which is ignored. Instead, a fraction
of the budget is reserved and distributed across every unit by irrigated area
before the remainder is targeted, and the cost of that choice is quantified
(Table~\ref{tab:equity}).

\begin{table}[h]\centering\small
\caption{The cost of an equity floor. A floor of 1.0 reproduces the baseline
exactly, which serves as an internal consistency check.}
\begin{tabular}{S[table-format=3.0]S[table-format=2.0]S[table-format=3.1]l}
\toprule
{Floor, \%} & {Units funded} & {Mm$^3$} & Cost vs.\ pure targeting \\
\midrule
0   & 7  & 273.0 & --- \\
25  & 18 & 264.2 & $-3.2\%$ \\
50  & 18 & 253.3 & $-7.2\%$ \\
75  & 18 & 237.9 & $-12.9\%$ \\
100 & 18 & 218.1 & $-20.1\%$ \\
\bottomrule
\end{tabular}
\label{tab:equity}
\end{table}

Reserving a quarter of the budget funds every governorate at a cost of
\SI{3.2}{\percent} of the water. This is offered as the practical
recommendation: it does not depoliticise the allocation, it prices the political
constraint and returns the dial to the committee.

\section{Two findings that contradicted our hypothesis}

\subsection{The rice belt is the most efficient land in Egypt, not the least}

The work was built expecting the northern Delta rice belt to lead, reasoning that
ponded paddy evaporates throughout the season. The measurement says the opposite.
Damietta returns an evaporated share of \num{0.157} and Kafr El Sheikh
\num{0.191}, against \num{0.23}--\num{0.29} elsewhere. A closed rice canopy
shades the water surface, so ponding does not imply evaporation.

\subsection{The ranking follows throughput, not practice}

\begin{table}[h]\centering\small
\caption{Area-weighted evaporated share by irrigation system.}
\begin{tabular}{lS[table-format=1.2]S[table-format=3.0]S[table-format=1.3]}
\toprule
System & {Irrigated, Mha} & {AETI, \si{\milli\metre}} & {$E/\text{AETI}$} \\
\midrule
Delta --- rice belt        & 1.76 & 592 & 0.210 \\
Nile Valley (Upper Egypt)  & 1.19 & 684 & 0.241 \\
Delta --- other            & 0.61 & 500 & 0.243 \\
\bottomrule
\end{tabular}
\end{table}

The wasted \emph{fraction} is close to uniform nationally. What differs is
throughput: Upper Egypt passes 660--\SI{775}{\milli\metre} through a field
against the Delta's 510--\SI{640}{\milli\metre}, so an equivalent inefficiency
costs more water. Targeting therefore follows evaporative demand rather than poor
practice --- a materially different policy message from the one we set out to
deliver.

\section{Self-audit: how much of this requires the satellite?}

The allocation was re-run with the satellite information progressively removed.

\begin{table}[h]\centering\small
\begin{tabular}{lS[table-format=2.1]}
\toprule
Information available & {Gain vs.\ baseline, \%} \\
\midrule
Real per-district evaporation (WaPOR)      & 25.2 \\
Per-system average only (agronomic map)    & 21.5 \\
Spatially uniform evaporation              & 0.0 \\
\bottomrule
\end{tabular}
\end{table}

With no spatial information, targeting gains nothing: the entire advantage is
spatial. But most of it resides in the system-level pattern, so the increment
attributable to district-level satellite \emph{detail} is approximately
\num{3.7} percentage points. This is a weaker claim than anticipated and is
reported because a reviewer would otherwise derive it. The counter-argument is
that the system-level pattern is precisely what the satellite corrected: without
it the budget would have been directed into the Delta on the strength of a
textbook intuition about ponded rice.

\section{Contamination we introduced and then removed}
\label{sec:contamination}

Aswan initially ranked first nationally. Inspection of the underlying pixel
distribution showed seasonal $E$ reaching \SI{1834}{\milli\metre}, with 5\,392
pixels above \SI{600}{\milli\metre} --- rates attainable by open water, not by
cropland. Lake Nasser, the Toshka lakes, the Nile channel, main canals and Delta
aquaculture had survived both the JRC mask, whose occurrence baseline misses
fluctuating margins and recent water bodies, and the beneficial-fraction floor,
which a partially vegetated shoreline pixel can clear.

The effect was not marginal: Aswan's mean $E$ was \SI{274}{\milli\metre} against a
median of \SI{192}{\milli\metre}. The mean was being set by contamination, and it
alone had placed Aswan at the head of the national ranking. Introducing the
physical ceiling $E < \SI{600}{\milli\metre}$ --- for every unit the 99th
percentile of cropland $E$ lies below \SI{400}{\milli\metre} while open water
exceeds \SI{1200}{\milli\metre}, so nothing real occupies the interval --- reduced
Aswan's mean to \SI{186}{\milli\metre} and largely removed its lead. The national
advantage fell from \SI{32.2}{\percent} to \SI{25.2}{\percent}, and the
system-level contrast that had appeared to separate Upper Egypt from the Delta
largely dissolved.

\keyfinding{This correction is reported because the uncorrected result would have
been more favourable to the hypothesis being tested. A method whose most
attractive output is produced by contamination is a method that has not yet been
checked.}

\section{Limitations}

\subsection{The principal methodological objection}

\citet{chukalla2022}, in the reference framework for WaPOR-based irrigation
performance assessment, \emph{excluded the beneficial fraction} on the grounds
that WaPOR's separation of evaporation from transpiration is insufficiently
accurate. That is a direct challenge to the present work, which rests entirely on
that separation, and it is not dismissed here.

Three considerations bear on it. First, the $E$/$T$ split is used
\emph{relatively}, to order units, rather than absolutely; a systematic bias in
the retrieval partly cancels under comparison between units processed by the same
algorithm. Second, the self-audit above shows the ordering is driven
substantially by AETI --- the best-constrained WaPOR variable --- through the
throughput effect, rather than by the fraction alone. Third, and decisively, the
appropriate status of this output is a \emph{prioritisation hypothesis to be
confirmed by field measurement}, not a settled allocation. The correct response
to \citet{chukalla2022} is to direct field validation at the units this method
ranks highest, which is itself a use of the ranking.

\subsection{Other limitations}

\begin{itemize}
\item \textbf{Field scale is not basin scale.} A share of the water an
  intervention frees was already returning to drains and aquifers and being
  reused downstream \citep{perry2007,vanhalsema2012}. Suppressing soil
  evaporation may also convert it to transpiration and yield rather than to
  recoverable water. Recoverable evaporation is an upper bound; these figures
  rank locations and do not constitute a water account.
\item \textbf{Magnitude is soft, ordering is comparatively robust.} The ordering
  follows observed evaporation; the size of the prize moves with every
  coefficient in Table~\ref{tab:interventions}.
\item \textbf{The decision unit is administrative.} Governorates are not the
  delivery unit; irrigation directorates and command areas are. The method is
  boundary-agnostic and runs unchanged on any polygon set, including those used
  by IPAT.
\item \textbf{Units below 5\,000\,ha are excluded} and reported, not
  silently dropped. GAUL's Luxor polygon covers 12\,km$^2$,
  which is the city rather than the governorate \citep{gaul2015}.
\item \textbf{The confidence indicator is presently unsound.} It is computed over
  the already-masked subset, so a predominantly desert unit can attain a narrow
  distribution and therefore a high score. It is reported in the tool and should
  not be relied upon.
\end{itemize}

\section{Implementation}

The analysis runs either as a local Python pipeline over exported GeoTIFFs or
entirely within Google Earth Engine \citep{gorelick2017}, in which case
per-unit statistics are precomputed to an Earth Engine asset and the allocation
is evaluated client-side. The delivered artefact is a single self-contained HTML
file requiring no network, which matters because the intended user opens it
inside a ministry rather than on a research network. All intervention
coefficients are editable at runtime, so the committee argues with the model
rather than being handed its conclusion.

Monitoring closes the loop. Because the indicator that justifies the spending is
the indicator that measures it, the following season observes funded and unfunded
units under the same instrument at the same cost, permitting a
difference-in-differences estimate with a stated confidence interval --- including
the case in which the effect is too small to detect.

\bibliographystyle{unsrtnat}
\bibliography{refs}

\end{document}
